\documentclass[11pt]{article}

\usepackage[margin=25mm]{geometry}

\usepackage{amsmath,amssymb,fontspec,unicode-math}

\IfFontExistsTF{DejaVu Serif}{\setmainfont{DejaVu Serif}}{\setmainfont{Latin Modern Roman}}

\IfFontExistsTF{Latin Modern Math}{\setmathfont{Latin Modern Math}}{\setmathfont{DejaVu Math TeX Gyre}}

\usepackage{xurl,hyperref}

\setlength{\parindent}{0pt}\setlength{\parskip}{6pt}

\begin{document}

\title{MODEM: A Morton-Order Degradation Estimation Mechanism for Adverse Weather Image Recovery}

\author{ELVIA preliminary AI analysis}\date{21 October 2026}\maketitle

Preliminary assessment: Minor.\par

Independent scientific review remains pending. Author adoption has not been requested.

The scan backend uses a first-order input coefficient instead of the displayed exact zero-order hold. Morton locality is a scope qualification; reported restoration scores remain untested.

\section*{Review scope}

Full supplied technical report read alongside the corresponding original-paper claim passages and applicable author-release source. This pass independently checks the stated reasoning and bounds; historical executions were not rerun.

\section*{Problem 1: Released scan uses an approximate input coefficient}

Source: Section 4.2, Eqs.(9)–(10), PDF p.7; modem\_arch.py, selective scan call; Mamba v2.0.4, selective\_scan\_ref

Assessment: supported; impact: minor

\subsection*{Source claim (paraphrase)}

The paper specifies exact zero-order hold \(\bar A=e^{\Delta A}\), \(\bar B=(e^{\Delta A}-I)A^{-1}B\). The current author release calls its pinned Mamba selective scan, whose input coefficient is \(\Delta B\).

\subsection*{Analysis}

For one state and step, take \(A=-1,B=C=1,\Delta=x_1=1,h_0=D=0\). The negative state parameter and positive step are allowed. The exact continuous equation \(\dot h=-h+1\) gives \(h(1)=1-e^{-1}\approx0.632121\), whereas the backend recurrence gives \(h_1=1\). Its transition coefficient is exponential, but its input increment uses no exponential correction.

\subsection*{Finding and scope}

The formula and linked primitive differ for the same generated coefficients. Both define legitimate discrete models; this does not refute restoration accuracy or imply that the exact-ZOH model would perform better.

\subsection*{Proposed correction}

Document the implemented first-order coefficient \(\Delta B\), or use the exact diagonal coefficient \(\operatorname{expm1}(\Delta a)b/a\), with limit \(\Delta b\) at zero. Changing a trained model’s recurrence requires new empirical validation.

The described correction remains a proposal; author adoption or a corresponding erratum is not established.

\section*{Problem 2: Morton ordering has block locality rather than a uniform neighbor guarantee}

Source: Appendix A.6, PDF p.19

Assessment: not\_issue; impact: minor

\subsection*{Source claim (paraphrase)}

The locality motivation says nearby image pixels remain nearby in the scan.

\subsection*{Analysis}

Adjacent pixels \((2^m-1,0)\) and \((2^m,0)\) have interleaved indices \((4^m-1)/3\) and \(4^m\), so their scan gap is \((2\cdot4^m+1)/3\) despite unit spatial distance. Conversely an aligned dyadic square fixes the high interleaved bits and occupies a contiguous interval.

\subsection*{Finding and scope}

The motivation is valid as hierarchical block locality and does not require a universal adjacent-pixel bound. No permutation or inverse-scan defect is established; the empirical scan comparison is not refuted.

\section*{Limits}

Preliminary AI source/argument assessment; independent human scientific adjudication remains pending.\par

No author-provided or retained verification program was executed. Published empirical results were not independently reproduced.\par

\section*{Source identity}

NeurIPS 2025 official proceedings PDF (exact retained local copy inspected)

\url{https://proceedings.neurips.cc/paper_files/paper/2025/file/43d2b7fbee8431f7cef0d0afed51c691-Paper-Conference.pdf}

Reviewed PDF SHA-256: \nolinkurl{c8e19cb6bc76c547f8009c178954d3bae0615e628810044552aa8f0ebc90ea09}

The SHA-256 identifies the exact retained PDF inspected in this rapid source/argument check. Original supplied report bytes are retained locally. No unambiguous original PDF retrieval timestamp was recovered; file modification time was not substituted.

\end{document}
